\section[Cote 156-6 : le lemme des images réciproques, et les ateliers triviaux]{Cote 156-6 : le lemme des images réciproques, et les ateliers triviaux\protect\verifie[\leanblobmain]{Folder156_6}}\label{sec:156-6}

La cote 156-6, chapitre~VI de \emph{Vers une géométrie des formes}, \q{Analysis situs (deuxième mouture)}, quatre-vingt-treize pages datées de sa main du 18 au 22~juin 1986, reprend huit jours après la première mouture (cote 156-4, section~\ref{sec:156-4}) l'axiomatique des \emph{ateliers} : un ensemble $\mathfrak F$ de figures muni de deux ordres, la sous-figure $\leq$ et le raffinement $\ll$, ses éléments irréductibles $\mathcal M$, les \emph{multistrates}, et les plus petites figures non vides pour $\ll$, les \emph{lieux}. Tout le dossier travaille avec les \emph{parties fermées} d'un ensemble ordonné, celles qui contiennent tout élément plus petit qu'un de leurs éléments ; on note $\mathfrak P_{\mathrm f}(I)$ leur ensemble. Ce sont les fermés de la topologie d'Alexandrov de l'ordre. On démontre ici trois énoncés des pages~40 à~50, qui ne portent que sur des ensembles ordonnés et leurs parties fermées ; les ateliers eux-mêmes (At\,1 à At\,7, les figures, $\ll$ sur $\mathfrak F$) ne sont pas codés.

\subsection*{Le lemme de la page 49}

Pour comparer deux ateliers, la page~48 veut lire une application entre parties fermées comme l'image réciproque par une application croissante en sens inverse, et la page~49 énonce : \q{Soient $I$, $I'$ deux ens.\ ordonnés. Alors l'application $\mathrm{Hom}_{\mathrm{croiss}}(I, I') \to \mathrm{Hom}_{\mathrm{cr}}(\mathfrak P_{\mathrm f}(I'), \mathfrak P_{\mathrm f}(I))$ [\dots] est \emph{injective}} ; son image est formée des applications qui commutent \q{aux intersections (quelconques)} et aux réunions (\lot{156-6}{3}{49}). La page est très retravaillée. Le mot \q{ordonnés} est écrit sur un \q{préordonnés} dont le \q{pré} est biffé, et la marge parle d'un \q{sobrifié} $\widehat I'$ de $I'$, dans des lignes qu'on ne lit que par fragments.

\begin{enonce}[Lemme, page 49]
Soient $I$ un ensemble préordonné et $I'$ un ensemble ordonné. L'application $\varphi \mapsto \varphi^{-1}$, des applications croissantes $I \to I'$ dans les applications $\mathfrak P_{\mathrm f}(I') \to \mathfrak P_{\mathrm f}(I)$, est injective. Son image est formée des applications qui commutent aux réunions et aux intersections quelconques. L'existence de $\varphi$ vaut encore si $I'$ n'est que préordonné.
\end{enonce}

\begin{proof}
Si $\varphi$ est croissante, l'image réciproque d'une partie fermée est fermée, et $\varphi^{-1}$ commute à toutes les réunions et intersections. Si $\varphi^{-1} = \psi^{-1}$, on applique cette égalité à la partie fermée $I'_{\leq \psi(x)}$, qui contient $\psi(x)$ : elle contient donc $\varphi(x)$, d'où $\varphi(x) \leq \psi(x)$, et de même $\psi(x) \leq \varphi(x)$. L'antisymétrie de l'ordre de $I'$ donne $\varphi = \psi$.

Inversement, soit $g$ une application qui commute aux réunions et intersections quelconques. Elle est croissante, puisque $A \subset B$ donne $g(A) = g(A \cap B) = g(A) \cap g(B)$. Pour $x \in I$, soit $A_x$ l'intersection des parties fermées $B$ telles que $x \in g(B)$ ; l'intersection vide est $I'$, et $g$ commute aux intersections, donc $x \in g(A_x)$. Or $A_x$ est la réunion des $I'_{\leq y}$, $y \in A_x$, et $g$ commute aux réunions : il existe $y \in A_x$ tel que $x \in g(I'_{\leq y})$, et l'on pose $\varphi(x) = y$. Alors $x \in g(B)$ équivaut à $\varphi(x) \in B$. Si $x \in g(B)$, $B$ contient $A_x$, donc $\varphi(x)$. Si $\varphi(x) \in B$, $B$ contient $I'_{\leq \varphi(x)}$, et $x \in g(I'_{\leq \varphi(x)}) \subset g(B)$. Ainsi $g = \varphi^{-1}$. Enfin $\varphi$ est croissante : si $x' \leq x$, alors $x$ est dans la partie fermée $g(I'_{\leq \varphi(x)})$, donc $x'$ aussi, c'est-à-dire $\varphi(x') \leq \varphi(x)$.
\end{proof}

\begin{encadre}
\begin{contrexemple}[Les deux hypothèses du lemme sont nécessaires]
\textup{(a)} Si $I'$ est seulement préordonné, l'injectivité tombe. Soit $I = I'$ formé de deux points dont chacun est $\leq$ à l'autre. Les seules parties fermées sont $\emptyset$ et $I'$, et l'identité et l'échange des deux points ont la même image réciproque.

\textup{(b)} On ne peut pas remplacer les réunions quelconques par les réunions finies. Soit $I'= \mathbf N$ avec son ordre usuel, et $I$ un point. L'application qui envoie la partie fermée $\mathbf N$ sur le point, et toute autre partie fermée sur $\emptyset$, commute aux intersections quelconques et aux réunions finies, mais elle n'est l'image réciproque par aucune application $\varphi$.
\end{contrexemple}
\end{encadre}

\begin{proof}
Pour (b) : les parties fermées de $\mathbf N$ sont $\emptyset$, les $\{0, \dots, n\}$ et $\mathbf N$ ; elles sont totalement ordonnées par inclusion, de sorte qu'une réunion de deux d'entre elles est l'une des deux, et une intersection vaut $\mathbf N$ si et seulement si chacun de ses termes vaut $\mathbf N$. Si l'application était $\varphi^{-1}$, avec $\varphi(\mathrm{pt}) = n$, la partie $\{0, \dots, n\}$ aurait pour image le point et serait donc égale à $\mathbf N$, ce qui est faux.
\end{proof}

\noindent Le contre-exemple~(a) justifie la correction de la page : le \q{pré} biffé devait l'être. Le contre-exemple~(b) éclaire la marge. Pour la topologie d'Alexandrov, $\mathbf N$ est un fermé irréductible, réunion croissante des $\{0, \dots, n\}$, sans point générique. L'application de (b) est l'image réciproque par le \q{point à l'infini} du sobrifié de $\mathbf N$, que la marge semble viser. La lecture fait le même rapprochement, et nomme l'énoncé exact du lemme correspondance de Raney (1952). Ces pages ne citent personne.

\subsection*{L'exemple de la page 40}

La page~39 appelle \emph{spécieuse} une partie fermée $\mathfrak F'$ de $\mathfrak F$ stable par les bornes supérieures des familles majorées \emph{dans $\mathfrak F$}. La page~40 prend $\mathfrak F = \mathfrak P(I)$, ordonné par inclusion. Donner une partie fermée revient à donner une \q{structure de \emph{complexe} simplicial sur $I$, du moins si $I$ est fini}, et \q{Dire que $\mathfrak F'$ est spécieux, signifie que la réunion de simplexes est un simplexe --- i.e.\ [\dots] donner une partie $I'$ de $I$, et $\mathfrak F' = \mathfrak P(I')$ --- c'est pas du tout pareil !} (\lot{156-6}{2}{40}). Dans $\mathfrak P(I)$ toute famille est majorée, par $I$.

\begin{enonce}[Page 40]
Soit $K \subset \mathfrak P(I)$ une famille stable par passage aux parties. Elle est stable par réunions quelconques si et seulement s'il existe $I' \subset I$ tel que $K = \mathfrak P(I')$. Si $I$ est fini et $K$ non vide, la stabilité par réunion de deux éléments suffit.
\end{enonce}

\begin{proof}
Si $K$ est stable par réunions, $I' = \bigcup K$ est dans $K$, et toute partie de $I'$ aussi ; la réciproque est claire. Si $K$ est non vide et stable par réunion de deux éléments, il contient $\emptyset$ et, par récurrence, toute partie finie de $\bigcup K$, puisque chaque singleton $\{a\}$, $a \in \bigcup K$, est partie d'un élément de $K$. Si $I$ est fini, $\bigcup K$ est fini, donc dans $K$, et l'on conclut comme plus haut.
\end{proof}

\noindent La restriction \q{du moins si $I$ est fini} est nécessaire pour la forme binaire. Les parties finies de $\mathbf N$ forment une famille non vide, stable par passage aux parties et par réunion de deux éléments, qui n'est aucun $\mathfrak P(I')$ : elle contiendrait $\mathbf N$, réunion de tous les singletons.

\subsection*{L'axiome des lieux dans les ateliers triviaux, pages 44 à 50}

La remarque de la page~44 fabrique un atelier à partir de n'importe quel ensemble ordonné $(\mathcal M, \leq)$, en posant $\ll\, = \,\leq$. Il ajoute : \q{dans les axiomes qui vont suivre, il y en a qui éliminent, je pense, cet exemple trivial d'atelier} (\lot{156-6}{3}{44}, \lot{156-6}{3}{45}). Dans un tel atelier, les lieux sont les éléments minimaux, et la petite ombre d'une multistrate $Y$ est $\mathrm{omb}(Y) = \{x \text{ minimal} \mid x \leq Y\}$. L'axiome des lieux At\,8 (page~45) demande que toute multistrate $X$ ait un lieu $x \in \mathrm{omb}(X)$ qui ne soit dans aucun $\mathrm{omb}(Y)$, $Y < X$. La page~49 le traduit pour les ateliers triviaux : \q{$\exists\, x \in \mathcal M$, $x$ \emph{minimal}, tel que $x \leq X$}, et qu'il y ait relation de \q{successeur} entre $x$ et $X$ \q{(dans le cas $X \neq x$ i.e.\ $X$ non minimal)}, la phrase passant à la page suivante : \q{i.e.\ $\nexists\, Y \in \mathcal M$ avec $x \leq Y < X$. C'est évidemment une condition des plus restrictives sur $\mathcal M$ !} (\lot{156-6}{3}{49}, \lot{156-6}{3}{50}).

\begin{enonce}[Pages 49 et 50, forme exacte]
Dans l'atelier trivial défini par un ensemble ordonné $(\mathcal M, \leq)$, l'axiome At\,8 a lieu si et seulement si tout élément de $\mathcal M$ est minimal, c'est-à-dire si l'ordre de $\mathcal M$ est l'égalité.
\end{enonce}

\begin{proof}
Si $X$ n'est pas minimal et si $x \in \mathrm{omb}(X)$, alors $x \neq X$, donc $x < X$, et $x$ est dans sa propre ombre $\mathrm{omb}(x)$ : c'est le cas $Y = x$ de la condition de la page. Aucun lieu ne convient. Si tous les éléments sont minimaux, on prend $x = X$, et il n'y a pas de $Y < X$. Enfin, si tous les éléments sont minimaux, $X \leq Y$ entraîne $Y \leq X$, donc $X = Y$.
\end{proof}

\noindent \q{Des plus restrictives} est donc en dessous de la vérité : At\,8 élimine tous les ateliers triviaux, sauf les discrets. La parenthèse \q{successeur} semble viser une autre condition, que $x$ soit couvert par $X$, sans $Y$ avec $x < Y < X$. Elle est strictement plus faible : la chaîne à deux éléments $0 < 1$ la satisfait, avec $x = 0$, mais ne satisfait pas At\,8. Le \q{$x \leq Y$} écrit sur la page est, lui, conforme à l'axiome de la page~45, et c'est lui qui force la discrétion. La lecture le relève en note, et Lean le confirme. Dans la première mouture, la propriété correspondante était une proposition, sans démonstration (cote 156-4, page~88) ; la deuxième mouture en fait un axiome.

Ne sont pas démontrés ici : la proposition de la page~41 (bijection entre parties fermées spécieuses de $\mathfrak F$ et parties fermées de $\mathcal M$, sous At\,1 et At\,2), le théorème-scholie des pages~32 et~33, et l'équivalence, annoncée à la page~49, entre la surjectivité de (1.54) et la condition sur les topologies ; aucun énoncé sur les ateliers eux-mêmes n'est codé. La seule piste ouverte pour cette cote dans le tableau des pistes porte sur les dates du dossier ; elle est codicologique, et la formalisation n'en dit rien.
